% META: {"id": "TSPE-LIMFCT-EX-030", "chapitre": "TSPE-LIMITES-FONCTIONS", "type_objet": "exercice", "capacites": ["TSPE-LIMFCT-C2"], "capacites_codes": ["C2"], "methodes": ["M2"], "parcours": 2, "competences": ["calculer", "raisonner"], "duree_min": 12, "mode_creation": "ex_nihilo", "sources_inspiration": [], "fichier_tex": "chapitres/TSPE-LIMITES-FONCTIONS/exercices/TSPE-LIMFCT-EX-030.tex", "corrige_tex": "chapitres/TSPE-LIMITES-FONCTIONS/corriges/TSPE-LIMFCT-CO-030.tex", "status": "generated"}
% BEGIN-VERIFY
% from sympy import *
% x = symbols('x')
% f = (2*exp(x) + 1)/(exp(x) + 3)
% assert limit(f, x, oo) == 2
% assert limit(f, x, -oo) == Rational(1, 3)
% END-VERIFY
\begin{exercice}{TSPE-LIMFCT-EX-030}{2}{12}
Soit $f(x) = \dfrac{2\mathrm{e}^x + 1}{\mathrm{e}^x + 3}$.
\begin{enumerate}
  \item Determiner les limites de $f$ en $+\infty$ et $-\infty$.
  \item En deduire les AH de $\mathcal{C}_f$.
  \item La courbe admet-elle une AV ?
\end{enumerate}
\end{exercice}
